\bookchapter{الفصل 1: مرحباً، حاسوبي الشخصي}

\chapterintro{هاتفك في الحقيقة هو حاسوبك الشخصي الفائق في جيبك. إنه أقوى بمرات لا تحصى من الحاسوب الذي وضع الإنسان على القمر، لكنه يرقد بهدوء في راحة يدك. اليوم، لنتعرف جيداً على هذا “العملاق الصامت”.}

\booksection{أولاً: “الأعضاء الداخلية” لعتاد الهاتف: تحليل معمّق}

الهاتف صغير لكن أعضاءه كاملة. دعونا نتعرف على هذه المكونات الدقيقة واحداً تلو الآخر كما في درس التشريح.

\booksubsection{SoC (شريحة النظام): “الدماغ + القلب” للهاتف}

SoC اختصار لـ System on Chip، وتعني “وضع النظام بأكمله على شريحة واحدة”. إنها ليست مجرد CPU (وحدة المعالجة المركزية)، بل تدمج CPU و GPU (وحدة معالجة الرسوميات) و NPU (وحدة معالجة الشبكات العصبية) و ISP (معالج إشارات الصورة) والنطاق الأساسي (المسؤول عن الاتصالات) وغيرها جميعاً معاً.

هذا مثل بناء مبنى الحكومة ومحطة الطاقة ومركز الاتصالات ومركز الشرطة في مدينة ما كلها في مبنى واحد، بكفاءة عالية جداً.

\booksubsection{CPU: جوهر الحوسبة العامة}

CPU هو المكوّن الأكثر مركزية في SoC، وهو مسؤول عن معالجة الغالبية العظمى من مهام الحوسبة. تستخدم معالجات الهاتف بنية ARM، والتي تختلف عن بنية x86 للحاسوب، وتتميز بنية ARM بمجموعة تعليمات مخفّضة (RISC)، واستهلاك طاقة أقل، وأنها أكثر ملاءمة للأجهزة المحمولة.

\textbf{تصميم أنوية CPU الكبيرة والصغيرة}: معالجات الهاتف الحديثة تعتمد بشكل شائع بنية “الأنوية الكبيرة والصغيرة” (تقنية big.LITTLE من ARM)، وتقسّم الأنوية إلى مستويات أداء مختلفة:

\begin{itemize}
  \item \textbf{الأنوية الفائقة/الكبيرة}: الأقوى أداءً، تتولى المهام الثقيلة. مثل تحميل الألعاب، تشغيل التطبيقات، الحوسبة المعقدة. التردد الأعلى (3GHz+)، لكن استهلاك الطاقة الأكبر أيضاً.
  \item \textbf{الأنوية المتوسطة}: توازن بين الأداء واستهلاك الطاقة، تعالج مهام الحمل المتوسط اليومية. تصفح Weibo ومشاهدة الفيديو والتبديل بين المهام يعتمد عليها أساساً.
  \item \textbf{الأنوية الصغيرة/الفعالة}: استهلاك الطاقة الأقل، تعالج مهام الخلفية الخفيفة. وضع الاستعداد، الاستماع للموسيقى، استقبال الرسائل، إلخ، تكاد لا تستهلك طاقة.
\end{itemize}

تعمل الأنوية الثلاثة المستويات معاً: يوزّع مجدول النظام المهام تلقائياً حسب الحمل — المهام الخفيفة للأنوية الصغيرة لتوفير الطاقة، والأعمال الشاقة للأنوية الكبيرة للعمل بقوة، ثم العودة للأنوية الصغيرة بعد الانتهاء. هذا هو السبب في أن الهاتف يمكنه تشغيل الألعاب بسلاسة والبقاء في وضع الاستعداد طوال اليوم.

\textbf{خصائص بنية كل علامة تجارية}:

\begin{itemize}
  \item \textbf{Qualcomm Snapdragon 8 Elite}: يعتمد بنية CPUمطورة ذاتياً Oryon (مصدرها استحواذ Nuvia)، نواة فائقة (4.32GHz) + 6 أنوية أداء (3.53GHz)، جميعها على عملية تصنيع TSMC الجديدة من الجيل الثاني بدقة 3nm. ودّع تماماً بنية ARM المرجعية، وحقق قفزة في الأداء أحادي النواة ومتعدد الأنوية. GPU هو Adreno 830 ببنية شرائح، تحسّن الأداء الرسومي بنسبة 40\%، وقوة حوسبة Hexagon NPU حوالي 45 TOPS، تحسن أداء استدلال AI بنسبة 45\%. (ملاحظة: نسخة Samsung Galaxy المخصصة لها تردد أعلى، النواة الفائقة تصل إلى 4.47GHz.)
  \item \textbf{MediaTek Dimensity 9400}: نواة فائقة (Cortex-X925، 3.62GHz) + 3 أنوية أداء (Cortex-X4، 3.3GHz) + 4 أنوية فعالة (Cortex-A720، 2.4GHz)، استمرار لبنية كل الأنوية الكبيرة، عملية تصنيع TSMC الجيل الثاني 3nm. GPU هو Immortalis-G925 (12 نواة)، تحسن الأداء الذروي بنسبة 41\%، تحسن أداء تتبع الأشعة بنسبة 40\%. معالج AI هو الجيل الثامن NPU 890، يدعم تشغيل نماذج AI كبيرة على الجهاز، وتدريب LoRA على الجهاز، وتوليد الفيديو على الجهاز.
  \item \textbf{Apple A18 Pro}: نواتا أداء (حوالي 4.04GHz) + 4 أنوية فعالة (حوالي 2.4GHz)، عملية تصنيع TSMC الجيل الثاني 3nm. GPU بـ 6 أنوية يدعم تتبع الأشعة على مستوى الأجهزة، NPU بـ 16 نواة (حوالي 35 TOPS) يوفر قوة حوسبة قوية لـ AI على الجهاز و Apple Intelligence. بنية Appleمطورة ذاتياً لا تزال متقدمة في الأداء أحادي النواة ونسبة الكفاءة، وبنية الذاكرة الموحدة توفر ميزة النطاق الترددي العالي لاستدلال AI.
  \item \textbf{HiSilicon Kirin 9020 من Huawei}: مدمج في سلسلة Mate 70، يعتمد بنية CPUمطورة ذاتياً Taishan، قدرات NPU و ISP تتعزز باستمرار، يدعم 5.5G والاتصال عبر الأقمار الصناعية، ويمثل أحدث إنجاز للرقائق المحلية.
\end{itemize}

\booksubsection{GPU: محرك عرض الرسوميات}

GPU (وحدة معالجة الرسوميات) مسؤول بشكل خاص عن مهام الرسوميات والحوسبة المتوازية. عدد أنويته أكبر بكثير من CPU (غالباً مئات أو آلاف الأنوية)، لكن وظيفة كل نواة أبسط نسبياً، مما يجعله مناسباً لمعالجة كميات كبيرة من الحسابات المتماثلة في وقت واحد.

\textbf{تدفق عرض الرسوميات}:

\begin{enumerate}
  \item \textbf{معالجة الرؤوس}: حساب إحداثيات موقع كل رأس في نموذج ثلاثي الأبعاد.
  \item \textbf{التنقيط}: تحويل الرسوميات ثلاثية الأبعاد إلى بكسلات على شاشة ثنائية الأبعاد.
  \item \textbf{تظليل البكسلات}: حساب اللون والملمس وتأثيرات الإضاءة لكل بكسل.
  \item \textbf{الإخراج والتركيب}: تركيب جميع الطبقات في الصورة النهائية وإرسالها للشاشة للعرض.
\end{enumerate}

\textbf{تأثير GPU على الاستخدام اليومي}: ليس فقط الألعاب! تمرير الواجهة، تأثيرات الرسوم المتحركة، تشغيل الفيديو، فلاتر الصور، تطبيقات الواقع المعزز كلها تعتمد على GPU. الهاتف ذو GPU الضعيف قد يسقط إطارات عند تمرير الواجهة، والألعاب تتوقف بشكل مباشر.

\textbf{GPU لكل علامة تجارية}: Qualcomm تستخدم GPUمطورة ذاتياً Adreno؛ MediaTek تستخدم ARM Mali GPU؛ Apple تستخدم GPUمطورة ذاتياً، والأداء والكفاءة كلاهما معيار في الصناعة.

\booksubsection{NPU: محرك تسريع AI}

NPU (وحدة المعالجة العصبية) معالج مصمم خصيصاً لعمليات الشبكات العصبية. مقارنة بـ CPU و GPU، فإن NPU أكثر كفاءة وأقل استهلاكاً للطاقة عند معالجة مهام AI.

\textbf{طريقة عمل الشبكات العصبية}: العملية الأساسية لـ NPU هي ضرب المصفوفات والعمليات التلافيفية — وهذا بالضبط ما هو أكثر الحسابات أساسية وتكراراً في الشبكات العصبية. من خلال التحسين على مستوى الأجهزة، يمكن لـ NPU معالجة كميات كبيرة من عمليات المصفوفات بالتوازي، بسرعة أكبر بعشرات المرات من CPU، واستهلاك طاقة أقل بكثير.

\textbf{تدفق معالجة AI}:

\begin{enumerate}
  \item \textbf{إدخال البيانات}: صور الكاميرا، صوت الميكروفون، بيانات المستشعرات، إلخ.
  \item \textbf{المعالجة المسبقة}: ضبط تنسيق البيانات وحجمها ليتناسب مع متطلبات إدخال الشبكة العصبية.
  \item \textbf{حساب الاستدلال}: يحمل NPU النموذج المدرّب، ويجري الحساب الأمامي على البيانات المدخلة، ويخرج النتيجة.
  \item \textbf{المعالجة اللاحقة}: تحويل مخرجات AI إلى نتائج قابلة للاستخدام، مثل إحداثيات الوجه، محتوى النص، علامات المشهد.
\end{enumerate}

\textbf{سيناريوهات التطبيق الفعلية}:

\begin{itemize}
  \item \textbf{التصوير الحسابي}: دمج المشاهد الليلية، تمويه الخلفية في الصور الشخصية، معالجة HDR، تجميل AI.
  \item \textbf{المساعد الصوتي}: التعرف على الصوت، فهم اللغة الطبيعية، تركيب الكلام.
  \item \textbf{الترجمة الفورية}: ترجمة AI للنصوص، ترجمة المحادثات.
  \item \textbf{التعرف الذكي}: التعرف على النصوص (OCR)، التعرف على الأشياء، التعرف على المشاهد.
  \item \textbf{النماذج الكبيرة على الجهاز}: أحدث الرقائق الرئيسية تدعم تشغيل نماذج AI صغيرة (مستوى 7B معامل) على الهاتف، لتحقيق حوار AI دون اتصال، وتلخيص المستندات، الكتابة الذكية وغيرها. قوة حوسبة Hexagon NPU في Snapdragon 8 Elite حوالي 45 TOPS، و NPU 890 في Dimensity 9400 يدعم أيضاً النماذج الكبيرة على الجهاز (سرعة استدلال LLM أسرع بنسبة 80\%)، و NPU بـ 16 نواة في Apple A18 Pro يوفر قوة حوسبة حوالي 35 TOPS.
  \item \textbf{معالجة الصور بالـ AI}: إزالة AI (إزالة المارة في الصور بنقرة واحدة)، توسيع الصور بالـ AI (توسيع حدود الصورة)، تحويل نمط AI، توليد الصور من النص، إلخ.
\end{itemize}

\booksubsection{ISP: معالج إشارات الصور}

ISP يعالج البيانات الخام القادمة من الكاميرا، ويحوّل إشارات الضوء إلى صور جميلة.

\textbf{كيف تقيس الأداء؟}انظر إلى الطراز! Apple A18 Pro و Qualcomm Snapdragon 8 Elite و MediaTek Dimensity 9400 و Huawei Kirin 9020 كلها مستوى رئيسي. المستوى المتوسط انظر إلى Snapdragon 8s Gen3/7+ Gen3 و Dimensity 8400. لا تنظر فقط إلى عدد الأنوية، فالبنية وعملية التصنيع (النانومتر) أهم. عملية 3nm هي الأكثر تقدماً حالياً، ونسبة الكفاءة أفضل بكثير من 5nm و 4nm.

\booksubsection{الذاكرة العشوائية (RAM): طاولة العمل المؤقتة}

الذاكرة العشوائية مثل مكتبك، كلما كبرت، زاد عدد الملفات التي يمكنك نشرها في وقت واحد. 12GB هي السائدة الآن، 16GB للمستخدمين الثقيلين، و 24GB فأكثر للاعبين والمتحمسين.

نوع الذاكرة العشوائية مهم أيضاً: LPDDR5 أسرع وأكثر توفيراً للطاقة من LPDDR4X. Apple رغم أن ذاكرتها ليست كبيرة، لكن تحسين النظام جيد، لذا تبدو “كافية”.

\textbf{ماذا يحدث إذا لم تكفِ الذاكرة؟}سيقوم النظام بإنهاء التطبيقات في الخلفية، وستحتاج إلى إعادة تحميلها عند التبديل. لهذا السبب أحياناً عند العودة من WeChat، تختفي اللعبة.

\booksubsection{التخزين (ROM/Flash): المستودع طويل الأمد}

التخزين مثل مكتبتك، المكان الذي توضع فيه جميع الصور والفيديوهات والتطبيقات. السائد الآن هو UFS 3.1 أو UFS 4.0، أسرع بمرات من eMMC القديم.

\textbf{لماذا سرعة التخزين مهمة؟}إنها تحدد سرعة فتح التطبيقات، سرعة تحميل الألعاب، سرعة حفظ الصور. بعد الاستخدام المطول، التخزين البطيء سيجعل الهاتف أبطأ فأبطأ.

\booksubsection{الشاشة: عيناك}

الشاشة هي النافذة الرئيسية للتفاعل بينك وبين الهاتف، وجودتها تؤثر مباشرة على التجربة.

\begin{itemize}
  \item \textbf{المادة}: OLED ألوان زاهية، تباين عالٍ، أكثر توفيراً للطاقة (عند عرض الأسود لا تضيء البكسلات)، لكن قد يكون هناك وميض عند السطوع المنخفض يؤذي العين؛ LCD أكثر راحة للعين، تكلفة أقل، لكن التباين والألوان لا تضاهي OLED.
  \item \textbf{الدقة}: 1080P هو السائد، 2K أوضح لكن يستهلك طاقة أكثر. مفهوم "شبكية العين" من Apple هو: على مسافة معينة، لا يمكن للعين البشرية تمييز بكسلات الشاشة.
  \item \textbf{معدل التحديث}: 60Hz قياسي، 90Hz/120Hz أكثر سلاسة، تمرير حريري. معدل التحديث العالي يستهلك طاقة، لذا العديد من الهواتف تدعم معدل التحديث التكيفي (LTPO)، الذي يتكيف ديناميكياً مع المحتوى.
  \item \textbf{معدل أخذ عينات اللمس}: يؤثر على سرعة استجابة اللمس، يجب أن يهتم به اللاعبون، كلما زاد كان أكثر استجابة.
  \item \textbf{طريقة التعتيم}: تعتيم DC أكثر راحة للعين من تعتيم PWM، عند السطوع المنخفض لا تتعب العين بسهولة. العديد من الهواتف الحديثة تدعم تعتيم PWM عالي التردد، أيضاً من أجل حماية العين.
\end{itemize}

\booksubsection{تقنية HDR للشاشة}

HDR (المدى الديناميكي العالي) يجعل الشاشة تعرض في وقت واحد أجزاء أكثر سطوعاً وأجزاء أكثر ظلمة، مما يرفع بشكل كبير إحساس الطبقات في الصورة.

\begin{itemize}
  \item \textbf{مبدأ التنفيذ}: نطاق سطوع الشاشة العادية محدود (SDR)، الفرق بين الأكثر سطوعاً والأكثر ظلمة ليس كبيراً بما يكفي. شاشة HDR لها سطوع ذروة أعلى (عادةً فوق 1000 شمعة)، تباين أكبر، نطاق لوني أوسع، يمكنها عرض تفاصيل أكثر.
  \item \textbf{المعايير السائدة}: HDR10 (الأكثر أساسية، مفتوح المصدر مجاني)، HDR10+ (بدعم من Samsung، بيانات وصفية ديناميكية)، Dolby Vision (رؤية Dolby، أفضل جودة لكن تتطلب ترخيصاً). المعيار الذي يدعمه الهاتف يعتمد على عتاد الشاشة وترخيص النظام.
  \item \textbf{التأثير الفعلي}: عند مشاهدة فيديو HDR، تكون أشعة الشمس أكثر سطوعاً، الظلال أعمق، الألوان أغنى. بشرط أن يكون الفيديو نفسه من إنتاج HDR، والشاشة تدعم عرض HDR.
\end{itemize}

\booksubsection{تقنية HDR للصور}

HDR في التصوير و HDR للشاشة أمران مختلفان. HDR للصور يحل عبر الخوارزميات مشكلة “التعريض الزائد في المناطق الساطعة والأسود التام في المناطق المظلمة”.

\begin{itemize}
  \item \textbf{مبدأ الخوارزمية}: عند الضغط على زر الغالق، يلتقط الهاتف فعلياً عدة صور بتعريضات مختلفة (تعريض ناقص للحفاظ على تفاصيل المناطق الساطعة، تعريض زائد للحفاظ على تفاصيل المناطق المظلمة، تعريض طبيعي للحفاظ على الدرجات المتوسطة)، ثم يدمجها في صورة واحدة بمدى ديناميكي عالٍ.
  \item \textbf{تدفق المعالجة}: التقاط متعدد الإطارات محاذاة وتعويض (مقاومة اهتزاز اليد) دمج التعريض/رسم الخرائط اللونية تصحيح الألوان إخراج الصورة النهائية.
  \item \textbf{تحسين التأثير}: الهواتف الحالية يمكنها بالفعل تحقيق "HDR في الوقت الفعلي" — ما تراه في عدسة الكاميرا هو تأثير معالجة HDR، ما تراه هو ما تحصل عليه. Smart HDR من Apple و HDR التلقائي من Android ناضجان جداً.
\end{itemize}

\booksubsection{البطارية والشحن: مصدر الطاقة}

البطارية هي أكثر المكوّنات “قصيرة العمر” في الهاتف، تتدهور بوضوح بعد سنتين أو ثلاث. السعة السائدة الآن 5500-7500mAh.

\begin{itemize}
  \item \textbf{قدرة الشحن}: معسكر Android وصل إلى 100W حتى 200W، يمتلئ في عشر دقائق؛ Apple لا تزال عند 20-30W. هل الشحن السريع يؤذي البطارية؟ مع الملحقات الأصلية لا مشكلة كبيرة، فالهاتف لديه حماية للتحكم في الحرارة.
  \item \textbf{الشحن اللاسلكي}: مريح لكن كفاءة منخفضة، حرارة عالية. مناسب للشحن العابر في المكتب، غير مناسب عند الاستعجال للخروج.
  \item \textbf{الشحن العكسي}: يتحول الهاتف إلى بنك طاقة، يشحن السماعات والساعة، للاستخدام الطارئ.
\end{itemize}

\booksubsection{الكاميرا: سحر الضوء والظل}

الكاميرات في الهواتف الحالية غالباً ما تكون ثلاثاً أو أربعاً، لكل منها تخصص:

\begin{itemize}
  \item \textbf{الكاميرا الرئيسية}: الأكثر استخداماً، المستشعر الأكبر، أفضل جودة للصورة.
  \item \textbf{الزاوية الواسعة جداً}: لالتقاط المناظر الطبيعية والمباني، مجال رؤية أوسع.
  \item \textbf{التقريب البعيد}: لالتقاط الأشياء البعيدة، هناك تقريب بصري وتقريب رقمي.
  \item \textbf{التصوير عن قرب}: لالتقاط الأشياء الصغيرة جداً عن مسافة قريبة.
\end{itemize}

\textbf{عدد البكسلات ليس كل شيء!}حجم المستشعر (المستشعر الأكبر يتفوق)، حجم فتحة العدسة (f/1.8 أكبر من f/2.4)، تثبيت الصورة (OIS التثبيت البصري مهم جداً) كلها تؤثر على جودة الصورة. التصوير الحسابي (خوارزميات AI) أصبح متزايد الأهمية، يمكنه دمج عدة صور في صورة واحدة جيدة.

\booksubsection{عدسة التقريب البعيد المنشورية}

عدسات التقريب البعيد التقليدية تحتاج إلى مسار بصري طويل، مما يجعل الهاتف سميكاً جداً. التقريب البعيد المنشوري عبر منشور بزاوية 90 درجة يكسر الضوء، مما يسمح للمسار البصري بالانتشار أفقياً داخل الهاتف، وبالتالي تحقيق تقريب بصري 5x أو حتى 10x مع الحفاظ على نحافة الهاتف.

\begin{itemize}
  \item \textbf{مبدأ البنية}: يدخل الضوء من العدسة يكسره المنشور 90 درجة يعبر مجموعة العدسات باتجاه عرض الهاتف يصل إلى المستشعر. يعادل وضع تلسكوب عمودي أفقياً.
  \item \textbf{الخصائص البصرية}: مضاعفات التقريب البصري عالية (5x-10x)، جودة الصورة تتفوق بكثير على التقريب الرقمي. لكن حجم المستشعر عادةً أصغر من الكاميرا الرئيسية، والأداء في الإضاءة المنخفضة متوسط.
  \item \textbf{ميزة التصوير}: تصوير القمر، تصوير المباني البعيدة، تصوير لقطات المسرح، هي ما لا يمكن للهواتف العادية بلوغه.
\end{itemize}

\booksubsection{OIS التثبيت البصري مقابل EIS التثبيت الإلكتروني}

تثبيت الصورة هو تقنية رئيسية تؤثر على جودة التصوير والفيديو، طريقتان بمبادئ مختلفة تماماً:

\begin{itemize}
  \item \textbf{OIS (التثبيت البصري)}: يعوض اهتزاز اليد من خلال التحريك الفيزيائي للعدسة أو المستشعر. داخل العدسة جيروسكوب مصغّر يكتشف اتجاه وحجم الاهتزاز، ثم يحرك المحرك لتحريك مجموعة العدسات في الاتجاه المعاكس. الميزة عدم فقدان جودة الصورة، العيب التكلفة العالية واحتلال المساحة.
  \item \textbf{EIS (التثبيت الإلكتروني)}: يعوض الاهتزاز من خلال خوارزمية برمجية تقطع وتزاحم الصورة. يسجل الهاتف نطاقاً أكبر قليلاً من الصورة، ثم يقطع المنطقة الثابتة بناءً على بيانات الجيروسكوب. الميزة عدم الحاجة إلى عتاد إضافي، العيب أنه يقطع الصورة (مجال الرؤية يصبح أصغر)، وقد تكون ضبابية في الإضاءة المنخفضة.
  \item \textbf{الاختيار الفعلي}: للتصوير الثابت أعطِ الأولوية لـ OIS، للفيديو الجمع بينهما يعطي أفضل تأثير. العديد من الهواتف الرئيسية الحالية تدعم تثبيت إزاحة المستشعر "Sensor-Shift"، أفضل من تثبيت العدسة التقليدي.
\end{itemize}

\booksubsection{التقريب البصري مقابل التقريب الرقمي}

\begin{itemize}
  \item \textbf{التقريب البصري}: يغير البعد البؤري من خلال التحريك الفيزيائي لمجموعة العدسات، دون فقدان جودة الصورة. التقريب البصري على الهاتف عادةً له مضاعفات ثابتة فقط (مثل 2x، 3x، 5x، 10x)، ولا يمكنه التكبير المستمر مثل الكاميرات.
  \item \textbf{التقريب الرقمي}: يكبّر المنطقة المركزية من الصورة، وهو في جوهره قطع + استيفاء تكبير. كلما زاد التكبير، زادت سوءاً جودة الصورة. التقريب الرقمي 100x يبدو رائعاً، لكن جودة الصورة الفعلية ضبابية جداً.
  \item \textbf{التقريب الهجين}: يجمع بين البصري والرقمي، يستخدم خوارزمية دمج متعدد الإطارات لتعويض فقدان الجودة بين مضاعفات التقريب البصري. التأثير أفضل من الرقمي البحت، لكن لا يزال أقل من البصري.
\end{itemize}

\booksubsection{الصوت: مكبرات الصوت وشريحة العرض المستقلة}

\textbf{مكبر صوت واحد مقابل مكبرين}:

\begin{itemize}
  \item \textbf{مكبر صوت واحد}: وحدة صوت واحدة فقط في الأسفل، الصوت يميل إلى جانب واحد، لا يوجد إحساس ثلاثي الأبعاد. شائع في الهواتف المتوسطة والمنخفضة.
  \item \textbf{مكبران (ستيريو)}: مكبر في الأعلى والأسفل، يشكلان مجال صوتي ثلاثي الأبعاد. عند مشاهدة الفيديو أو الاستماع للموسيقى، قنوات الصوت اليمنى واليسرى واضحة، الإحساس بالانغماس أقوى بكثير. استخدام السماعة كمكبر علوي هو حل شائع.
  \item \textbf{اختلاف جودة الصوت}: المكبرات المزدوجة ليست مجرد "صوت أعلى"، المفتاح هو الإحساس المكاني والإحساس بالطبقات. مشاهدة فيلم بمكبر واحد تجعلك تشعر أن شيئاً ما مفقود.
\end{itemize}

\textbf{شريحة العرض المستقلة في هواتف الألعاب}:

بعض هواتف الألعاب والهواتف عالية الأداء مزودة بشريحة معالجة عرض مستقلة، ووظيفتها:

\begin{itemize}
  \item \textbf{إدراج الإطارات}: إدراج إطارات لعبة 30fps إلى 60fps أو 120fps عبر الخوارزمية، الصورة أكثر سلاسة. تتم هذه العملية على شريحة العرض المستقلة، دون استهلاك موارد GPU في SoC.
  \item \textbf{الترقية الفائقة}: ترقية الصور منخفضة الدقة إلى دقة عالية عبر خوارزميات AI، جودة الصورة أوضح.
  \item \textbf{تقليل استهلاك الطاقة}: لأن إدراج الإطارات والترقية الفائقة تقوم بها شريحة العرض المستقلة، ينخفض حمل SoC، وتقل الحرارة واستهلاك الطاقة، اللعبة أكثر استقراراً.
\end{itemize}

\booksubsection{تحديد المواقع GPS: أحادي التردد وثنائي التردد}

يعتمد تحديد مواقع الهاتف أساساً على أنظمة الملاحة بالأقمار الصناعية مثل GPS.

\begin{itemize}
  \item \textbf{GPS أحادي التردد}: يستقبل إشارة نطاق L1 فقط، دقة تحديد المواقع حوالي 3-5 أمتار. كافٍ في المناطق المفتوحة، لكن في أودية المدن ذات المباني الشاهقة، انعكاس الإشارة يسبب انحرافاً كبيراً في تحديد المواقع.
  \item \textbf{GPS ثنائي التردد}: يستقبل إشارتي نطاقي L1 و L5 في وقت واحد. نطاق L5 له طول موجي أطول، قدرة أقوى على مقاومة التداخل، والحل المزدوج التردد يمكن أن يلغي معظم أخطاء الغلاف الأيوني، دقة تحديد المواقع تصل إلى 1-3 أمتار. ميزة واضحة في الملاحة في المدينة وتحديد المواقع لطلب السيارات.
  \item \textbf{دعم أبراج متعددة}: الهواتف الحديثة لا تدعم GPS فقط، بل تدعم أيضاً BeiDou (الصين)، GLONASS (روسيا)، Galileo (أوروبا) وغيرها من أنظمة الأقمار الصناعية. كلما زاد عدد الأقمار الصناعية المتصلة في وقت واحد، كان تحديد المواقع أسرع وأدق.
\end{itemize}

\booksubsection{المستشعرات: قرون استشعار الهاتف الصغيرة}

يختبئ في الهاتف أكثر من عشرة مستشعرات، تعمل بصمت:

\begin{itemize}
  \item \textbf{مقياس التسارع + الجيروسكوب}: يحسان بوضعية الهاتف، الانعطاف في الألعاب، التدوير التلقائي للشاشة.
  \item \textbf{مستشعر الضوء}: يضبط سطوع الشاشة تلقائياً.
  \item \textbf{مستشعر المسافة}: إطفاء الشاشة تلقائياً عند مكالمة هاتفية وقرب الهاتف من الأذن.
  \item \textbf{مستشعر البصمة}: السائد الآن بصمة تحت الشاشة ضوئية أو فوق صوتية.
  \item \textbf{التعرف على الوجه}: Face ID من Apple هو ضوء منظم ثلاثي الأبعاد، مستوى أمان عالٍ؛ العديد من هواتف Android هي تعرف ثنائي الأبعاد بالصور، الأمان متوسط.
  \item \textbf{قفل Bootloader}: اتصال قريب المدى، تمرير بطاقة الحافلة، بطاقة الدخول، البطاقة البنكية.
  \item \textbf{البارومتر}: قياس الارتفاع عن سطح البحر، يكتشف حتى صعود الدرج.
  \item \textbf{مستشعر هول}: يكتشف فتح وإغلاق الغطاء المغناطيسي.
\end{itemize}

\booksection{ثانياً: نظام التشغيل: Android مقابل iOS مقابل HarmonyOS}

العتاد هو الجسد، ونظام التشغيل هو الروح. لكل من الأنظمة الثلاثة خصائصها:

\begin{itemize}
  \item \textbf{أندرويد}: مثل مجتمع مفتوح المصدر، يمكن للشركات المصنعة للهواتف تخصيصه بعمق، خيارات كثيرة. يمكنك تغيير السمات، تثبيت متاجر تطبيقات خارجية، تجربة أشياء مختلفة. لكن بسبب الانفتاح، التجزئة خطيرة، واختلاف التجربة بين العلامات التجارية كبير.
  \item \textbf{iOS}: مثل حديقة منمّقة بعناية، Apple تتحكم في العتاد والبرمجيات معاً، التجربة موحدة وسلسة. جودة التطبيقات عالية، حماية الخصوصية جيدة، لكن درجة الحرية منخفضة، الكثير من الأشياء يجب أن تتم وفق قواعد Apple.
  \item \textbf{HarmonyOS}: كائن جديد من Huawei، يركز على القدرات الموزعة، يمكنه جعل الأجهزة المختلفة تتعاون بسلاسة. سنتحدث عنه بالتفصيل لاحقاً.
\end{itemize}

الأنظمة الثلاثة كلها غنية بالنظام البيئي، اختر أسلوب الحياة الذي تحبه. لا يوجد جيد أو سيئ مطلق، فقط ما يناسبك أو لا.

\booksection{ثالثاً: تحليل شامل لمنافذ الهاتف}

منافذ الهاتف هي الجسر الذي يربطه بالعالم الخارجي، والمنافذ المختلفة تختلف كثيراً:

\booksubsection{Micro USB}

\begin{itemize}
  \item \textbf{الخصائص الفيزيائية}: منفذ شبه منحرف مسطح، يميز الوجه الأمامي والخلفي (لا يدخل إذا أُدخل بالعكس). تم التخلص منه أساساً في الهواتف، لكن بعض الأجهزة منخفضة المستوى والقديمة لا تزال تستخدمه.
  \item \textbf{قيود النقل}: يدعم فقط معيار USB 2.0، أقصى سرعة نقل 480Mbps، قدرة الشحن عادةً لا تتجاوز 10W.
  \item \textbf{العيوب}: الوجهان غير قابلين للتبديل، عمر قصير للإدخال والإخراج، سرعة نقل بطيئة، قدرة شحن منخفضة.
\end{itemize}

\booksubsection{Lightning (منفذ البرق من Apple، تم التخلص منه)}

\begin{itemize}
  \item \textbf{الخصائص الفيزيائية}: منفذ 8 دبابيسمطورة ذاتياً من Apple، يمكن إدخاله من الوجهين. يُستخدم في iPhone 14 والإصدارات الأقدم، وقد تم التخلص منه بالكامل الآن.
  \item \textbf{قيود النقل}: منفذ Lightning فيزيائياً يدعم فقط معيار USB 2.0، أقصى سرعة نقل 480Mbps. حتى مع استخدام كابل USB-C إلى Lightning، فإن سرعة النقل محدودة بمنفذ Lightning نفسه.
  \item \textbf{سبب التخلص منه}: الاتحاد الأوروبي أصدر تشريعاً يلزم جميع الهواتف باستخدام منفذ USB-C بشكل موحد بدءاً من 2024، و Apple بدأت من سلسلة iPhone 15 بالتحول بالكامل إلى Type-C. منفذ Lightning خرج رسمياً من مسرح التاريخ.
  \item \textbf{ملاحظة}: كابلات البيانات المعتمدة من Apple MFi تحتوي على شريحة مصادقة، الكابلات غير المعتمدة قد لا تشحن أو تنقل البيانات بشكل طبيعي. إذا كنت لا تزال تستخدم iPhone قديم، يُنصح بتجهيز كابل بيانات Lightning.
\end{itemize}

\booksubsection{USB Type-C}

\begin{itemize}
  \item \textbf{الخصائص الفيزيائية}: منفذ بيضاوي الشكل بـ 24 دبوس، يمكن إدخاله من الوجهين. أصبح المنفذ القياسي للهواتف ومعظم الأجهزة الإلكترونية. الاتحاد الأوروبي أصدر تشريعاً يلزم جميع الهواتف والأجهزة اللوحية وسماعات الرأس وغيرها من الأجهزة المحمولة باستخدام منفذ USB-C بشكل موحد بدءاً من 2024، ومناطق أخرى في العالم تتبع.
  \item \textbf{بروتوكول النقل}: Type-C هو شكل المنفذ فقط، قدرة النقل الفعلية تعتمد على البروتوكول الداخلي. يدعم USB 2.0 (480Mbps)، USB 3.0/3.1/3.2 (5-20Gbps) وحتى USB4 (40Gbps).
  \item \textbf{تمييز مهم}: ليس كل منفذ Type-C يدعم USB 3.0! كثير من الهواتف تستخدم منفذ Type-C لكنها متصل داخلياً بمسامير USB 2.0 فقط، فسرعة النقل بطيئة كالمنافذ القديمة. إصدار iPhone 16/17 القياسي يدعم سرعة USB 2.0 فقط، أما إصدار Pro فيدعم USB 3.0 (يتطلب كابل USB-C to USB-C).
  \item \textbf{قدرة الشحن}: يدعم بروتوكولات الشحن السريع مثل USB PD، يمكن أن يدعم حتى 240W شحن.
  \item \textbf{إخراج الفيديو}: يدعم DisplayPort Alt Mode، يمكن إخراج الفيديو إلى شاشة عبر كابل تحويل.
\end{itemize}

\booksubsection{منفذ صوت 3.5mm}

\begin{itemize}
  \item \textbf{الخصائص الفيزيائية}: فتحة دائرية بـ 3 أو 4 مقاطع، 3 مقاطع تنقل صوت قناة يمين ويسار، 4 مقاطع تضيف قناة ميكروفون.
  \item \textbf{الميزات}: نقل إشارة تناظرية، تشغيل فوري، لا يحتاج إلى تعريفات أو تحويل. جودة الصوت مستقرة، التأخير صفر.
  \item \textbf{الوضع الحالي}: معظم الهواتف الرئيسية ألغت منفذ 3.5mm، تحتاج إلى كابل تحويل Type-C أو استخدام سماعات بلوتوث. بعض الهواتف المتوسطة وهواتف الألعاب لا تزال تحتفظ به.
\end{itemize}

\booksection{رابعاً: نظام منتجات الهاتف وقواعد التسمية}

في مواجهة تشكيلة المنتجات المبهرة من كل علامة تجارية، فإن فهم قواعد التسمية يمكن أن يساعدك على تحديد الهاتف المناسب لك بسرعة.

\booksubsection{سلاسل منتجات العلامات التجارية الرئيسية:}

\begin{itemize}
  \item \textbf{Apple}: iPhone (الإصدار القياسي) / iPhone Plus (إصدار الشاشة الكبيرة) / iPhone Pro (الإصدار الاحترافي) / iPhone Pro Max (إصدار احترافي بشاشة كبيرة) / iPhone SE (الإصدار الاقتصادي) / iPhone 16e (إصدار ذكاء اصطناعي اقتصادي، مزود بشريحة A18 ومرسام C1 من تطوير Apple). في 2026، سلسلة iPhone 17 مزودة بالكامل بشريحة A19 Pro، مع تحسين قدرات الذكاء الاصطناعي.
  \item \textbf{Huawei}: سلسلة Mate (رائد للأعمال) / سلسلة P/Pura (رائد للتصوير) / سلسلة Mate X (رائد قابل للطي) / سلسلة nova (شباب متوسط) / سلسلة Enjoy (مبتدئ)
  \item \textbf{Xiaomi}: سلسلة Xiaomi الرقمية (رائد) / Xiaomi MIX (رائد استكشافي/قابل للطي) / سلسلة Redmi K (رائد للأداء) / سلسلة Redmi Note (متوسط) / سلسلة Redmi الرقمية (مبتدئ)
  \item \textbf{OPPO}: سلسلة Find X (رائد) / سلسلة Reno (متوسط) / سلسلة A (مبتدئ) / سلسلة OnePlus الرقمية (رائد للأداء)
  \item \textbf{vivo}: سلسلة X (رائد للتصوير) / سلسلة iQOO الرقمية (رائد للأداء) / سلسلة S (سيلفي متوسط) / سلسلة Y (مبتدئ)
  \item \textbf{Samsung}: سلسلة Galaxy S (رائد) / Galaxy Z Fold/Flip (رائد قابل للطي) / سلسلة Galaxy A (متوسط)
\end{itemize}

\booksubsection{معاني اللواحق الشائعة:}

\begin{itemize}
  \item \textbf{Pro}: اختصار Professional، الإصدار الاحترافي. عادةً يرقّي الكاميرا والشاشة والمعالج وغيرها من التكوينات الأساسية على أساس الإصدار القياسي.
  \item \textbf{Max}: الإصدار الأكبر. الشاشة الأكبر، البطارية الأكبر، لكن التكوين ليس بالضرورة أعلى من Pro.
  \item \textbf{Ultra}: الإصدار الفائق. عادةً يمثل أعلى تكوين في السلسلة، المواصفات الأقوى، السعر الأعلى أيضاً. مثل Samsung S24 Ultra، Xiaomi 14 Ultra.
  \item \textbf{SE}: Special Edition الإصدار الخاص، في الواقع غالباً الإصدار الاقتصادي. التكوين مخفّض، السعر أقل. مثل iPhone SE.
  \item \textbf{Plus}: الإصدار المكبّر. الشاشة أكبر من الإصدار القياسي، البطارية أكبر أيضاً، لكن التكوين أساسي متماثل.
  \item \textbf{Lite/إصدار الشباب}: الإصدار الخفيف. التكوين مبسّط، السعر أقل.
  \item \textbf{Note}: كان يشير في الأصل إلى هاتف بشاشة كبيرة، الآن يختلف استخدامه بين العلامات التجارية. Redmi Note سلسلة متوسطة.
  \item \textbf{Neo}: الإصدار الجديد. عادةً يشير إلى هاتف متوسط بأداء عالٍ.
\end{itemize}

\booksection{خامساً: النطاق الأساسي والهوائيات: الأبطال الخفيون لاتصالات الهاتف}

الهاتف يمكنه الاتصال والتصفح بفضل ثنائي شريحة النطاق الأساسي والهوائي.

\booksubsection{شريحة النطاق الأساسي:}

النطاق الأساسي هو المعالج المسؤول عن الاتصالات اللاسلكية في الهاتف، فهو يطبّق الإشارات الرقمية إلى موجات راديو ويخرجها، أيضاً يفكّ ترميز موجات الراديو المستلمة إلى إشارات رقمية. أداء النطاق الأساسي يؤثر مباشرة على جودة المكالمات، سرعة الإنترنت، استقرار الإشارة. Qualcomm هي أكبر مورّد لشرائح النطاق الأساسي للهواتف، MediaTek و Samsung لديهما أيضاً شرائح نطاق أساسيمطورة ذاتياً. Apple بدأت من iPhone 16e بحمل شريحة النطاق الأساسي C1مطورة ذاتياً، مما يمثل دخول Apple رسمياً عصر شرائح النطاق الأساسيمطورة ذاتياً، شريحة C1 تستخدم عملية تصنيع 4nm، أداء استهلاك الطاقة ممتاز؛ iPhone 16/16 Plus/16 Pro/16 Pro Max لا تزال تستخدم شرائح نطاق أساسي من Qualcomm (Snapdragon X75 5G)، ومن المتوقع أن تنتقل الإصدارات اللاحقة تدريجياً إلى شرائحمطورة ذاتياً. الإصدارات الأقدم iPhone XS/11 التي كانت تستخدم شرائح Intel كان أداء الإشارة فيها ضعيفاً، وكان نقطة ألم لكثير من المستخدمين.

\booksubsection{تصميم الهوائي:}

هوائيات الهاتف مختبئة داخل الجسم، عادةً مرتبة حول الإطار. تصميم الهوائي يؤثر مباشرة على أداء استقبال الإشارة:

\begin{itemize}
  \item \textbf{تقنية MIMO}: إدخال متعدد وإخراج متعدد، استخدام هوائيات متعددة لإرسال واستقبال البيانات في وقت واحد، تحسين السرعة والاستقرار. 4×4 MIMO أسرع من 2×2 MIMO.
  \item \textbf{عدد الهوائيات}: الهواتف الرئيسية عادةً لها أكثر من عشرة هوائيات (شبكة الخلوي، Wi-Fi، بلوتوث، NFC، GPS كل منها يحتاج هوائي مستقل).
  \item \textbf{تأثير الجسم المعدني}: المعدن يحمي الإشارة، لذا الهواتف ذات الجسم المعدني تحتاج إلى تصميم شريط هوائي (شريط بلاستيكي على الإطار). الأجسام الزجاجية والبلاستيكية لها تأثير أقل على الإشارة.
\end{itemize}

\booksubsection{تطور تقنيات الشبكة:}

\begin{itemize}
  \item \textbf{4G LTE}: سرعة الذروة النظرية تنزيل 1Gbps/رفع 500Mbps (Cat.18)، التجربة الفعلية عادةً 50-150Mbps. التغطية الأوسع، كافية.
  \item \textbf{5G NR}: سرعة الذروة النظرية تنزيل 10Gbps+. نطاق Sub-6GHz تغطية أفضل، mmWave الموجة المليمترية سرعة فائقة لكن تغطية صغيرة جداً.
  \item \textbf{5G-A (5G Advanced)}: النسخة المحسّنة من 5G، وتسمى أيضاً 5.5G. سرعة التنزيل تصل إلى 10Gbps والرفع 1Gbps، مع تحسن كبير في كثافة الاتصالات ودقة التموضع. منذ 2024 بدأت China Mobile وChina Telecom وChina Unicom نشر شبكات 5G-A في المدن الكبرى، وهواتف Huawei Mate 70/80 وXiaomi 15/16 وغيرها من الهواتف الرائدة تدعم 5G-A. كما يتيح 5G-A قدرة "التواصل والإحساس الموحّد"، حيث يمكن للمحطة الاتصال والإحساس بالبيئة في آن واحد، مما يوفر البنية التحتية للاقتصاد منخفض الارتفاع (توصيل الطائرات المسيّرة) وإنترنت المركبات.
\end{itemize}

\booksubsection{تصنيف نطاقات الشبكة:}

النطاقات التي يدعمها الهاتف تحدد في أي دول وعلى شبكات أي مشغلين يمكن استخدامه:

\begin{itemize}
  \item \textbf{FDD-LTE Band (سلسلة B)}: مثل B1، B3، B5، إلخ، نطاقات 4G. دول مختلفة تستخدم نطاقات مختلفة، الصين تستخدم أساساً B1/B3/B5/B8/B34/B38/B39/B40/B41.
  \item \textbf{TDD-LTE Band}: معيار آخر لـ 4G، الصين تستخدم بكثرة B38/B39/B40/B41.
  \item \textbf{5G NR Band (سلسلة N)}: مثل n1، n3، n28، n41، n78، n79، إلخ. n78 (3.5GHz) هو النطاق الرئيسي لـ 5G في الصين، n41 (2.5GHz) يستخدم بكثرة من قبل China Mobile.
  \item \textbf{نصيحة الشراء}: إذا كنت بحاجة لاستخدامه في الخارج، فتأكد قبل الشراء من أن الهاتف يدعم نطاقات الدولة المستهدفة. الهواتف الصينية عادةً تدعم النطاقات الصينية بشكل كامل، لكن قد تفتقر إلى بعض النطاقات الخارجية.
\end{itemize}

\booksection{سادساً: هاتفك، كم دوراً يلعب في الواقع؟}

لم يعد مجرد هاتف:

\begin{enumerate}
  \item \textbf{أداة اتصال}: الهاتف، الرسائل النصية، WeChat، البريد الإلكتروني...
  \item \textbf{كاميرا ذكية}: التصوير، الفيديو، مسح المستندات...
  \item \textbf{محفظة}: الدفع، بطاقة الحافلة، البطاقة البنكية، بطاقة العضوية...
  \item \textbf{بوابة المعلومات}: الأخبار، الطقس، الأسهم، المعرفة...
  \item \textbf{جهاز ألعاب}: من Match-3 إلى ألعاب AAA...
  \item \textbf{مكتب متنقل}: المستندات، الجداول، الاجتماعات، التعاون...
  \item \textbf{مساعد صحي}: عد الخطوات، معدل ضربات القلب، النوم، الرياضة...
  \item \textbf{جهاز تحكم المنزل الذكي}: التحكم بالأضواء، المكيف، التلفاز، الروبوت المنظف...
  \item \textbf{هوية رقمية}: بطاقة الهوية الإلكترونية، رخصة القيادة الإلكترونية، رمز الصحة...
  \item \textbf{مركز ترفيهي}: الموسيقى، الفيديو، الكتب الإلكترونية، البودكاست...
\end{enumerate}

فكر، جهاز صغير يحمل الكثير من الحياة! لذا، تعلم استخدامه جيداً هو تعلم العيش بشكل أفضل.

\begin{tcbpractice}{لنطبق: تعرف على تكوين العتاد}
في [الإعدادات]، انظر إن كان يمكنك إيجاد طراز المعالج، حجم الذاكرة العشوائية (RAM) والتخزين الداخلي (ROM). كن على دراية بقدرات هاتفك. مستخدمو Android يمكنهم تنزيل “AnTuTu” أو “DevCheck” لعرض معلومات العتاد التفصيلية.

\end{tcbpractice}
